\section{Conclusion}
\label{sec:conclusion}

\subsection{Summary}

Sequential Monte Carlo is the correct tool for statistical inference about the
space of valid redistricting plans.
Its three defining properties are:

\paragraph{Calibrated.}
The importance weights correct for the spanning-tree proposal's non-uniformity
by construction (Proposition~\ref{prop:weight-correct}).
No Markov chain mixing assumption is required.
For any inferential claim of the form ``the enacted plan is at the $X$th
percentile of all valid plans by partisan outcome,'' only \smc{} provides a
statistically valid answer from a feasibly sized sample.

\paragraph{Parallelisable.}
Particle proposals at each stage are independent.
Rayon \texttt{par\_iter()} gives approximately linear speedup in the number
of cores, making \smc{} with $N = 5{,}000$ particles practical on an 8-core
workstation within 3--8 minutes for mid-size states.

\paragraph{Auditable.}
The \texttt{file\_sha256} binding and per-resample \texttt{index\_map} records
allow a verifier to trace any particle's complete construction history from
\texttt{base\_seed} alone, with stronger guarantees than seed-only recording.

\subsection{Deployment Guidance}

The standard workflow for a redistricting practitioner is:
\begin{center}
  \cs{} (generate plan) $\to$ \smc{} (calibrated ensemble) $\to$
  percentile claim (court filing).
\end{center}

Concretely:
\begin{enumerate}
  \item Run \cs{} to generate the submitted plan from the
        minimum-edge-cut seed over $T = 600$ METIS seeds.
  \item Run \smc{} with $N = 5{,}000$ particles (for $k = 8$--$14$) to
        generate a calibrated ensemble of valid plans.
  \item Compute the percentile of the submitted plan's partisan seat outcome
        in the weighted \smc{} ensemble.
  \item Report the percentile, the \ess{}, and the \texttt{base\_seed} in
        the legislative findings.
  \item Store the full NDJSON output file with its \texttt{file\_sha256} for
        verifier access.
\end{enumerate}

\subsection{Open Questions}

\begin{enumerate}
  \item \textbf{Validation against R \texttt{redist}.}
        The reference implementation is \texttt{redist::redist\_smc()} in
        the R \texttt{redist} package~\citep{mccartan2023}.
        Phase 2 will run both implementations on the same census-tract graph
        and base seed and compare the marginal distributions of partisan seat
        outcomes, confirming that the Rust implementation is algorithmically
        equivalent.

  \item \textbf{SeedCompositor::SmcPercentile.}
        The current CLI supports ensemble generation but not plan selection
        from the weighted ensemble.
        Phase 2 will add \texttt{SeedCompositor::SmcPercentile}, enabling
        \texttt{redist ensemble --method smc --select-percentile 0.5} to
        return the plan at the weighted median of the ensemble.

  \item \textbf{Adaptive particle count.}
        The particle count $N$ is currently fixed before the run.
        An adaptive scheme that increases $N$ when the \ess{} falls below
        a target would provide automatic quality control without user-set
        parameters.

  \item \textbf{Large-$k$ scaling.}
        For California ($k = 52$) and Texas ($k = 38$), the ESS degradation
        model (Proposition~\ref{prop:ess-degrade}) predicts that $N = 20{,}000$
        particles will be needed to maintain $\ess \geq 10{,}000$.
        This has not been tested empirically; the memory requirement ($\approx
        2$ GB for CA) is within reach but has not been validated.

  \item \textbf{Connection to Polsby-Popper weighting.}
        The current implementation targets the uniform distribution over
        valid plans.
        A natural extension is to target a distribution that upweights compact
        plans (e.g., by the Polsby-Popper score), which would produce an
        ensemble representative of ``geographically realistic'' plans rather
        than the fully uniform distribution.
        This requires modifying the importance weights to include the
        compactness ratio.
\end{enumerate}

\subsection{Position in the Toolkit}

The redistricting toolkit is now complete at all five method levels.
Table~\ref{tab:comparison} (Section~\ref{sec:comparison}) summarises when
to use each method.
The full workflow for evidentiary redistricting is:
\begin{center}
  \cs{} (certify plan quality) $\to$ \flip{} (local sensitivity)
  $\to$ \be{} (conditional structure audit)
  $\to$ \smc{} (global calibrated inference).
\end{center}
Each layer adds a different type of evidence.
\smc{} is the final layer: the one that converts the ensemble into a
statistically valid court-facing percentile claim.
